% =====================================================================
%  CV — version générale, candidatures recherche (stages, labos, LinkedIn)
%  Dérivé de google-student-researcher.tex par SUPPRESSION du ciblage.
%
%  CE QUI A ÉTÉ RETIRÉ par rapport à la version Google :
%   - "Available from 02/27 for a six-month placement in Canada" → TODO(human)
%   - "see attached transcript" → aucun relevé n'accompagne un CV envoyé à froid
%   - pdftitle spécifique à l'annonce
%
%  CE QUI EST RESTÉ, et pourquoi :
%   - Publications AVANT l'expérience : vrai pour toute candidature recherche,
%     pas seulement pour Google. C'est le premier critère d'un encadrant.
%   - Diplôme en MM/YY : c'était une exigence Google, mais un format de date
%     non ambigu à l'international ne coûte rien et évite le 03/04 vs 04/03.
%
%  INTERDITS (master-cv.md §11) : projet clical, −58 %, PFA 82,3 %,
%  audience Medium, Google Final Round, projets perso, noms de clients.
% =====================================================================

\documentclass[10pt, letterpaper]{article}

\usepackage[
    ignoreheadfoot,
    top=1.1 cm,
    bottom=1.1 cm,
    left=1.2 cm,
    right=1.2 cm,
    footskip=1.0 cm,
]{geometry}
\usepackage{titlesec}
\usepackage{enumitem}
\usepackage{fontawesome5}
\usepackage[dvipsnames]{xcolor}
\usepackage{textcomp}

\definecolor{MyRed}{RGB}{160, 20, 40}

\usepackage[
    pdftitle={Melissa Colin — CV — Research},
    pdfauthor={Melissa Colin},
    colorlinks=true,
    urlcolor=MyRed
]{hyperref}

\usepackage{charter}

\pagestyle{empty}
\setlength{\parindent}{0pt}

\titleformat{\section}
  {\large\bfseries\scshape\raggedright\color{MyRed}}{}{0em}{}[\titlerule]
\titlespacing{\section}{0pt}{7pt}{3pt}

\setlist[itemize]{leftmargin=*, noitemsep, topsep=1pt, partopsep=0pt, parsep=0pt}

\newcommand{\entry}[4]{%
  \noindent \textbf{#1} \hfill #2 \\
  \textit{#3} \hfill \textit{#4} \par
}

\IfFileExists{personal.tex}{\input{personal.tex}}{}
\providecommand{\myphone}{+33 7 82 52 XX XX}

\begin{document}

% ---------------------------------------------------------------- HEADER
\begin{center}
    {\fontsize{21}{21}\selectfont \textbf{Mélissa COLIN}} \\[4pt]

    % Positionnement : large d'abord, spécialité en second. Le lecteur voit
    % qu'il peut te confier autre chose, sans que tu perdes ta signature.
    % Le statut d'étudiante a été retiré de cette ligne : la section Education
    % arrive juste dessous et la date "Expected 09/27" le dit déjà. Ça libère
    % la place pour la spécialité, qui, elle, n'est écrite nulle part ailleurs.
    {\large Deep Learning \& Computer Vision $\cdot$ Specialising in 3D Human Motion} \\[4pt]

    \small{
    \faMapMarker* Bordeaux, France \quad|\quad
    % TODO(human) — LE NUMÉRO DE TÉLÉPHONE.
    % Sur un CV envoyé par mail à un encadrant : garde-le, c'est un canal direct.
    % Sur un PDF déposé dans la section "Sélection" de LinkedIn : ce fichier est
    % téléchargeable par n'importe qui, indexable, et tu ne peux plus le reprendre.
    % Décide. Si tu publies sur LinkedIn, commente la ligne ci-dessous.
    \faPhone* \myphone{} \quad|\quad
    \faEnvelope~\href{mailto:contact@melissacolin.ai}{contact@melissacolin.ai} \\
    \faGlobe~\href{https://melissacolin.ai}{melissacolin.ai} \quad|\quad
    \faLinkedin~\href{https://www.linkedin.com/in/melissacolin/}{LinkedIn} \quad|\quad
    \faGithub~\href{https://github.com/melissa-colin}{github.com/melissa-colin} \quad|\quad
    \faGraduationCap~\href{https://scholar.google.com/citations?user=7r7iFpsAAAAJ}{Scholar} \quad|\quad
    \href{https://orcid.org/0009-0003-2525-4824}{ORCID 0009-0003-2525-4824}
    }
\end{center}

\vspace{2pt}

% ---------------------------------------------------------------- ACCROCHE
% Phrases 1 et 2 : ce que tu fais et par quel mécanisme. Indépendant de l'annonce.
% Phrase 3 : "seeking", pas "followed by". "Followed by" se lit comme un plan
% déjà arrangé, donc comme une candidate déjà prise : un encadrant qui cherche
% un doctorant passe au dossier suivant. "Seeking" transforme la même phrase en
% invitation. "academic or industrial" ouvre aussi la porte aux labos privés.
% Phrase 4 : fenêtre bornée des deux côtés. Un encadrant a besoin du début ET
% de la fin pour savoir si son projet tient dedans, sinon il doit t'écrire.
{\small \textit{3D human motion reconstruction and computer vision. I design neural
architectures for multi-person motion and the differentiable physics engine they train
against, so that physically impossible motion falls outside what the model can express.
Expected graduation 09/27, and seeking a PhD position in computer vision afterwards,
academic or industrial. Available from 02/27 for a research placement of up to eight months,
through 09/27.}}

% ---------------------------------------------------------------- EDUCATION
\section{Education}

\entry{ENSEIRB-MATMECA, Bordeaux INP}{Sep 2024 -- Expected 09/27}%
{MEng in Computer Science (\textit{Diplôme d'ingénieur}), AI track, currently in final year}{Bordeaux, France}
\begin{itemize}
    % "see attached transcript" retiré : ce CV circule seul. Un lecteur étranger
    % qui ne connaît pas le /20 doit quand même pouvoir situer le chiffre, d'où
    % le rang et le pourcentage qui, eux, se lisent dans n'importe quel système.
    \item Rank in cohort: \textbf{5\textsuperscript{th}/81 (top 6\%)} in semester 8, Feb--May 2026, and \textbf{3\textsuperscript{rd}/93} in semester 6; 4\textsuperscript{th}/81 on the project unit, 5\textsuperscript{th}/53 on the advanced electives. French /20 scale; transcript available on request.
    \item Feb--May 2026: \textbf{Deep Learning 18/20} $\cdot$ \textbf{Research Initiation 17/20} (research-methods elective) $\cdot$ Applied TCP/IP 17.5/20 $\cdot$ Cryptology 16.5/20.
    \item Coursework: deep learning, computability \& complexity, game theory, operations research, numerical algorithms, compilation.
\end{itemize}

\vspace{4pt}

\entry{EPSI Bordeaux}{Sep 2021 -- Sep 2024}%
{BSc in Artificial Intelligence \& Data Science (RNCP level 6, EQF/bachelor level)}{Bordeaux, France}
\begin{itemize}
    \item Professional dossier across four competency blocks: data modelling, predictive modelling, production \& maintenance of AI systems, project management.
\end{itemize}

% ---------------------------------------------------------------- PUBLICATIONS
\section{Publications}

\begin{itemize}
    \item \textbf{Colin, M.}, Wang, Y., Cheng, L. \textit{InterForce: Physically Admissible Multi-Person Motion Refinement.} First author, in preparation, targeting \textbf{CVPR 2027}.
    \item \textbf{Colin, M.}, Chraibi Kaadoud, I. \textit{Performances and Explainability of ViT and CNN Architectures: An Empirical Study Using LIME, SHAP and Grad-CAM.} RJCIA / PFIA 2024, La Rochelle. First author, published, presented orally.
    \hfill \href{https://hal.science/hal-04641791}{[HAL]}
    % Titre exact du papier de Yilin Wang encore inconnu : marqué [À CONFIRMER]
    % dans master-cv.md §3.2, donc je ne l'écris pas. Demande-le AVANT le
    % 11 septembre : tu perds l'accès au Slack UAlberta en partant.
    \item Wang, Y., Li, S., \textbf{Colin, M.}, et al. \textit{X-Dancers: Group Choreography as Masked Completion of a Formation-Grounded Motion Tensor.} Submitted to \textbf{ICLR 2027}. Contribution: physically refined the input and output motion data, and produced the visualizations.
\end{itemize}

% ---------------------------------------------------------------- RESEARCH
\section{Research Experience}

\entry{Vision and Learning Lab, University of Alberta (Prof. Li Cheng)}{May 2026 -- Sep 2026}%
{Visiting Research Student, \textbf{Mitacs Globalink Research Award (2026)}}{Edmonton, Canada}
\begin{itemize}
    \item \textbf{Evidence-based motion correction.} Built a pipeline that corrects reconstructed 3D human motion \textit{without reading the noisy motion itself}, using only signals extracted from the image: ray--plane foot-ground contact, z-buffer occlusion ordering, monocular depth, camera calibration. The correction therefore comes from a source independent of the error.
    \item \textbf{Depth corrector.} Monocular depth from \textit{Depth Anything}, brought to a usable scale by anchoring it on two quantities measurable directly in the image (foot height and head size), then used to fix the depth ordering between people. Purely algorithmic: nothing trained, no ground truth. Result: $-$33\% inter-person penetration, measured on 100\% of the data.
    \item \textbf{Model design} (first-author work, \textit{InterForce}): a permutation-equivariant transformer whose attention is factorised over people, joints and time, so a scene never fixes its number of people; and an output in \textbf{forces, not poses}, backed by a \textbf{differentiable rigid-body dynamics engine written from scratch} (recursive Newton--Euler, composite-rigid-body mass matrix, joint limits) inside the training loop --- an impossible motion becomes unrepresentable rather than penalised. Clips a physics simulator replays without the character falling: \textbf{19\% $\to$ 79\%}.
    \item \textbf{27{,}000-clip paired corpus} built via \textit{render-then-estimate} (clean motion rendered to video, then re-estimated), processed on a 12-shard multi-GPU pipeline with work stealing across two compute sites.
\end{itemize}

\vspace{4pt}

\entry{ENSEIRB-MATMECA, Research Initiation elective, 17/20}{Feb -- May 2026}%
{Reproduction and critical study of \textit{InterMask} (ICLR 2025)}{Bordeaux, France}
\begin{itemize}
    \item Reproduced the VQ-VAE baseline (FID 1.034 $\to$ \textbf{0.918}), then tested two hypotheses on the physical constraints of the loss (asymmetric foot-contact penalty, temporal augmentation).
    \item \textbf{Reported a negative result:} neither variant improved FID or R-Precision. Published the finding as-is rather than dropping it.
\end{itemize}

% ---------------------------------------------------------------- INDUSTRY
\section{Industry Experience}

\entry{Sector Group}{Jun -- Aug 2025}{R\&D Engineer, AI and Safety}{Paris, France}
\begin{itemize}
    \item Trustworthy-AI research on information extraction from \textbf{non-OCR'd PDF documents} (no text layer, so vision and document processing had to be combined), under the \textbf{IEC 61508} functional-safety framework. RAG pipeline (RAGFlow fork), 100+ queries, 9 business units audited.
    % ATTENTION : formulation "abstention" à vérifier avant envoi. master-cv.md
    % §4.3 marque le couple 70 % / 100 % comme [À CONFIRMER]. Ce couple n'est
    % cohérent que si le système s'abstient quand il n'est pas sûr : 100 % de
    % précision sur les réponses émises, et le recall bas en est la
    % contrepartie logique. Si la sémantique réelle est autre, réécris la ligne.
    \item Tuned the pipeline for an abstention regime, where a wrong answer costs more than a missing one: \textbf{100\% precision on the answers actually returned}, at 70\% recall over the evaluation query set.
\end{itemize}

\vspace{4pt}

\entry{Cali Intelligences (deeptech start-up, real-time video analytics)}{Jan 2023 -- Sep 2024}%
{Applied AI Intern, then AI Developer (apprenticeship)}{Talence, France}
\begin{itemize}
    \item Designed and built the full training pipeline for the detection models (YOLOv7 $\to$ YOLOv8): ingestion, label conversion, augmentation library written from scratch in OpenCV, pre-training validation, YAML-driven orchestration, Docker/GitHub Actions/K3s deployment. \textbf{Detection accuracy 60\% $\to$ 90\%.}
    \item Benchmarked and selected the annotation tooling against 38 explicit requirements; redesigned the data storage into a bronze/silver medallion architecture on AWS for GDPR and EU AI Act traceability.
\end{itemize}

% ---------------------------------------------------------------- SKILLS
\section{Technical Skills}
\begin{itemize}
    \item \textbf{Programming languages:} \textbf{Python} (advanced) $\cdot$ C, SQL (intermediate) $\cdot$ Java $\cdot$ TypeScript/JavaScript $\cdot$ C++ $\cdot$ Bash $\cdot$ R.
    \item \textbf{Deep learning:} PyTorch, CNNs, Transformers/ViT, VQ-VAE, LoRA/PEFT, multi-stage fine-tuning. Implemented from scratch: K-Means, Naive Bayes, logistic regression, backprop through a hidden layer, PCA.
    \item \textbf{3D human motion:} SMPL/SMPL-H/SMPL-X, axis-angle rotations, forward kinematics, AMASS, monocular 3D pose estimation, Multi-HMR, camera calibration, z-buffer rasterisation. Metrics: MPJPE, PA-MPJPE, penetration, foot-skate, jitter, PFC, FID-motion, R-Precision, SDF.
    \item \textbf{Explainability:} LIME, SHAP, Grad-CAM, Grad-CAM++ (subject of my RJCIA 2024 paper).
    \item \textbf{Tooling:} Docker, Kubernetes/K3s, Git, GitHub Actions, Linux, SLURM-style GPU clusters, \LaTeX{}.
    \item \textbf{Languages:} French (native) $\cdot$ English \textbf{B2, TOEIC 840/990}, research internship entirely in English $\cdot$ Mandarin Chinese (elementary).
\end{itemize}

% ---------------------------------------------------------------- AWARDS
\section{Awards \& Leadership}
\begin{itemize}
    \item \textbf{Mitacs Globalink Research Award (2026)}: competitive international research funding, awarded on project proposal; funds the University of Alberta placement.
    \item \textbf{Vice-President, Forum INGENIB} (Apr 2025 -- present): \texteuro{}100k budget, 20-student team, 800 students and 80 companies hosted.
    \item \textbf{Networking Lead, ai4industry workshop} (Oct 2024 -- Jun 2025): recruited and coordinated 22 speakers (PhD students and researchers) around post-doctorate career panels.
    \item Student engagement formally validated at \textit{"high level"} by the school for two consecutive semesters (16/20).
\end{itemize}

\end{document}
